\section{Array: Array Terindeks dan Asosiatif, Fungsi Array Dasar}

Array adalah struktur data yang menyimpan banyak nilai dalam satu variabel. Bayangkan array sebagai sebuah lemari dengan banyak laci: setiap laci memiliki nomor (indeks) atau label (kunci) dan menyimpan satu nilai. Array sangat penting dalam pemrograman web karena data dari form, database, dan API sering kali berbentuk kumpulan nilai yang perlu diproses secara bersamaan. PHP mendukung dua jenis array utama: \textbf{array terindeks} (indeks numerik) dan \textbf{array asosiatif} (kunci string) \cite{phpManual}.

\textbf{Array terindeks} menggunakan indeks numerik yang dimulai dari 0. Anda dapat membuat array terindeks dengan sintaks pendek \code{[...]} atau fungsi \code{array(...)}. Elemen dapat ditambahkan dengan \code{\$arr[] = nilai;} (indeks otomatis) atau \code{\$arr[5] = nilai;} (indeks spesifik). \textbf{Array asosiatif} menggunakan string sebagai kunci, sangat berguna untuk merepresentasikan data terstruktur seperti profil pengguna atau konfigurasi. PHP juga mendukung array bersarang (\textit{nested array}) --- array di dalam array --- untuk merepresentasikan data multi-dimensi \cite{w3schoolsPhp}.

\begin{webcode}[language=PHP, caption={Array terindeks, asosiatif, dan bersarang}]
<?php
// Array terindeks
$buah = ["Apel", "Mangga", "Jeruk", "Durian"];
echo $buah[0];   // Apel
echo count($buah); // 4

// Array asosiatif
$mhs = [
    "nama" => "Rina",
    "nim"  => "2024001",
    "ipk"  => 3.75
];
echo $mhs["nama"];  // Rina

// Array bersarang (multi-dimensi)
$kelas = [
    ["Andi", 85],
    ["Budi", 92],
    ["Cici", 78]
];
echo $kelas[1][0];  // Budi
echo $kelas[1][1];  // 92

// Menambah elemen
$buah[] = "Semangka";       // indeks otomatis (4)
$mhs["email"] = "rina@mail.com";  // kunci baru
?>
\end{webcode}

PHP menyediakan puluhan fungsi bawaan untuk memanipulasi array. Berikut adalah fungsi-fungsi array yang paling sering digunakan.

\begin{table}[ht]
\centering
\begin{tabular}{|P{3.5cm}|P{4cm}|P{4.5cm}|}
\hline
\textbf{Fungsi} & \textbf{Kegunaan} & \textbf{Contoh Hasil} \\
\hline
\code{count(\$arr)} & Jumlah elemen & \code{count([1,2,3])} $\rightarrow$ \code{3} \\
\hline
\code{array\_push(\$a, v)} & Tambah di akhir & Menambah elemen baru \\
\hline
\code{array\_pop(\$a)} & Hapus dari akhir & Mengembalikan elemen terakhir \\
\hline
\code{array\_shift(\$a)} & Hapus dari awal & Mengembalikan elemen pertama \\
\hline
\code{in\_array(v, \$a)} & Cek keberadaan & Mengembalikan \code{true}/\code{false} \\
\hline
\code{array\_merge(a, b)} & Gabung dua array & Array baru yang digabung \\
\hline
\code{array\_keys(\$a)} & Ambil semua kunci & Array berisi kunci \\
\hline
\code{array\_values(\$a)} & Ambil semua nilai & Array berisi nilai \\
\hline
\code{sort(\$a)} & Urutkan naik & Mengubah array langsung \\
\hline
\code{array\_map(f, \$a)} & Terapkan fungsi & Array baru hasil transformasi \\
\hline
\end{tabular}
\caption{Fungsi array PHP yang sering digunakan}
\end{table}

Iterasi array paling idiomatik menggunakan \code{foreach}, seperti yang telah dibahas di section perulangan. Untuk memproses setiap elemen dan menghasilkan array baru, fungsi \code{array\_map()} sangat berguna karena menerapkan fungsi callback ke setiap elemen. Fungsi \code{array\_filter()} menyaring elemen berdasarkan kondisi, sedangkan \code{array\_reduce()} mengakumulasikan seluruh elemen menjadi satu nilai. Ketiga fungsi ini merupakan pola pemrograman fungsional yang semakin populer di PHP modern \cite{phpRightWay}.

\begin{catatan}
\textbf{Array di PHP Bukan Array Sebenarnya:} Secara internal, array PHP adalah \textit{ordered hash map} (peta hash terurut). Ini berarti array PHP dapat sekaligus memiliki kunci numerik \textbf{dan} string, serta urutan elemen selalu terjaga. Fleksibilitas ini membuat array PHP sangat powerful tetapi berbeda dari array di bahasa lain seperti Java atau C.
\end{catatan}

